% ============================================================================
% APPENDIX PA: LIT-AGHION - Schumpeterian Growth & Creative Destruction
% ============================================================================
% Category: Literature Integration (LIT-)
% Language: English
% EBF Integration: Innovation, growth theory, and institutional dynamics
% ============================================================================

\chapter*{Appendix UNMAPPED_PA: Philippe Aghion Research Integration}
\addcontentsline{toc}{chapter}{Appendix UNMAPPED_PA: LIT-AGHION - Schumpeterian Growth \& Creative Destruction}
\label{app:aghion}

% ============================================================================
% HEADER BLOCK
% ============================================================================
\begin{tcolorbox}[colback=blue!5!white, colframe=blue!75!black, title=Appendix Metadata]
\textbf{Code:} PA \\
\textbf{Category:} LIT- (Literature Integration) \\
\textbf{Title:} Schumpeterian Growth \& Creative Destruction \\
\textbf{Version:} 1.0 \\
\textbf{Status:} Complete \\
\textbf{Primary Author Focus:} Philippe Aghion \\
\textbf{Papers Covered:} 30 \\
\textbf{Dependencies:} DOMAIN-MACRO, DOMAIN-IO, LIT-ACEMOGLU (L), CONTEXT-CULTURE \\
\textbf{EBF Connections:} Innovation dynamics, institutional complementarities, long-run welfare
\end{tcolorbox}

% ============================================================================
% CROSS-REFERENCE MAP
% ============================================================================
\begin{tcolorbox}[colback=green!5!white, colframe=green!75!black, title=Cross-Reference Map]
\textbf{Outgoing References:}
\begin{itemize}[nosep]
    \item DOMAIN-MACRO: Growth theory foundations
    \item DOMAIN-IO: Competition-innovation relationship
    \item DOMAIN-SUSTAINABILITY: Green innovation and carbon policy
    \item LIT-ACEMOGLU (L): Institutions and development
    \item LIT-BLOOM (Q): Firm dynamics and innovation
    \item CONTEXT-CULTURE: Trust, regulation, and institutions
\end{itemize}

\textbf{Incoming References:}
\begin{itemize}[nosep]
    \item Chapter 9 (Context): Institutional Ψ parameters
    \item Chapter 14 (Applications): Innovation policy design
    \item PREDICT-05: Long-run growth predictions
\end{itemize}
\end{tcolorbox}

% ============================================================================
% ABSTRACT
% ============================================================================
\section*{Abstract}

Philippe Aghion has revolutionized growth economics through his development of \textbf{Schumpeterian growth theory}, which places innovation and creative destruction at the center of economic dynamics. His work provides crucial foundations for understanding how competition, institutions, and policy shape long-run prosperity. This appendix integrates Aghion's 30 key papers into the EBF framework, with particular emphasis on:
\begin{enumerate}
    \item The \textbf{Inverted-U relationship} between competition and innovation
    \item \textbf{Creative destruction} and its welfare implications
    \item \textbf{Directed technical change} and green innovation
    \item \textbf{Institutions, trust, and regulation}
    \item \textbf{AI, automation, and the future of growth}
\end{enumerate}

% ============================================================================
% QUICK REFERENCE
% ============================================================================
\begin{tcolorbox}[colback=yellow!10!white, colframe=orange!75!black, title=Quick Reference: Key Aghion Parameters]
\begin{tabular}{lll}
\textbf{Parameter} & \textbf{Value} & \textbf{Source} \\
\hline
Competition-innovation peak & HHI $\approx$ 0.15 & Aghion et al. (2005) \\
Creative destruction rate & 4-6\% annually & Aghion \& Howitt (1992) \\
Missing growth from CD & 0.5-1.0 pp/year & Aghion et al. (2019) \\
Carbon tax innovation elasticity & 0.3-0.5 & Aghion et al. (2016) \\
Trust-regulation tradeoff & $\rho = -0.4$ & Aghion et al. (2010) \\
\end{tabular}
\end{tcolorbox}

% ============================================================================
% FUNDAMENTAL QUESTION
% ============================================================================
% -----------------------------------------------------------------------------
% APPENDIX SCOPE BOX
% -----------------------------------------------------------------------------

\begin{tcolorbox}[
    colback=orange!5!white,
    colframe=orange!75!black,
    title={\textbf{Appendix Scope: Ziel / In-Scope / Out-of-Scope / Constraints}},
    fonttitle=\bfseries
]

\textbf{Ziel (Objective):}
\begin{itemize}[nosep]
    \item What drives long-run economic growth, and how do competition, institutions, and policy shape the innovation process that underlies prosperity?
\end{itemize}

\textbf{In-Scope:}
\begin{itemize}[nosep]
    \item Fundamental Question
    \item Creative Destruction and Welfare
    \item Institutions, Trust, and Regulation
    \item Directed Technical Change and Green Innovation
    \item Artificial Intelligence and the Future of Growth
\end{itemize}

\textbf{Out-of-Scope (delegiert an andere Appendices):}
\begin{itemize}[nosep]
    \item Glossary $\rightarrow$ Appendix G
\end{itemize}

\textbf{Constraints (Anwendungsgrenzen):}
\begin{itemize}[nosep]
    \item [TODO: Define when this appendix does NOT apply]
    \item [TODO: Define required prerequisites]
    \item [TODO: Define known limitations]
\end{itemize}

\textbf{Lieferobjekte (Deliverables):}
\begin{itemize}[nosep]
    \item 1 Table(s)
    \item 6 Equation(s)
    \item Worked Example(s)
\end{itemize}

\end{tcolorbox}


\section{Fundamental Question}

\begin{tcolorbox}[colback=red!5!white, colframe=red!75!black, title=Central Question]
\textbf{What drives long-run economic growth, and how do competition, institutions, and policy shape the innovation process that underlies prosperity?}
\end{tcolorbox}

Aghion's answer: Growth is driven by \textbf{quality-improving innovations} that create temporary monopoly rents but also \textbf{destroy} existing technologies. The optimal policy environment balances incentives for innovation with competitive pressure for efficiency.

% ============================================================================
% THEORETICAL FOUNDATIONS
% ============================================================================
\section{Theoretical Foundations: Schumpeterian Growth}

\subsection{The Basic Model}

Aghion \& Howitt (1992) formalized Schumpeter's vision of creative destruction:

\begin{definition}[Schumpeterian Growth]
Economic growth occurs through a sequence of quality-improving innovations. Each innovation:
\begin{enumerate}
    \item Renders the previous technology obsolete (creative destruction)
    \item Grants the innovator temporary monopoly power
    \item Generates positive externalities for future innovation
\end{enumerate}
\end{definition}

\begin{axiom}[UNMAPPED_PA-1: Creative Destruction]
\label{ax:pa1}
The arrival rate of innovations $\lambda$ determines both growth and obsolescence:
\begin{equation}
g = \lambda \cdot \ln(\gamma)
\end{equation}
where $\gamma > 1$ is the step size of quality improvements. Higher $\lambda$ increases growth but also increases the rate at which incumbents are replaced.
\textbf{EBF Connection:} Provides dynamic foundations for DOMAIN-MACRO welfare calculations.
\end{axiom}

\subsection{Competition and Innovation: The Inverted-U}

Aghion et al. (2005) discovered the non-monotonic relationship between competition and innovation:

\begin{theorem}[Inverted-U Relationship]
The relationship between product market competition (PMC) and innovation follows an inverted-U:
\begin{equation}
Innovation = \alpha_0 + \alpha_1 \cdot PMC - \alpha_2 \cdot PMC^2
\end{equation}
with the peak occurring at intermediate competition levels.
\end{theorem}

\begin{axiom}[UNMAPPED_PA-2: Escape Competition Effect]
\label{ax:pa2}
In neck-and-neck industries (similar technology levels), competition increases the incentive to innovate to ``escape'' to monopoly:
\begin{equation}
\frac{\partial \text{Innovation}}{\partial \text{Competition}} > 0 \quad \text{when } \Delta_{tech} \approx 0
\end{equation}
\textbf{EBF Connection:} Enriches DOMAIN-IO with context-dependent competition effects.
\end{axiom}

\begin{axiom}[UNMAPPED_PA-3: Schumpeterian Effect]
\label{ax:pa3}
In industries with large technology gaps, competition reduces innovation incentives by shrinking post-innovation rents:
\begin{equation}
\frac{\partial \text{Innovation}}{\partial \text{Competition}} < 0 \quad \text{when } \Delta_{tech} >> 0
\end{equation}
\textbf{EBF Connection:} Explains why ``one-size-fits-all'' competition policy fails.
\end{axiom}

% ============================================================================
% CREATIVE DESTRUCTION AND WELFARE
% ============================================================================
\section{Creative Destruction and Welfare}

\subsection{The Welfare Implications}

Aghion et al. (2016) examine how creative destruction affects subjective well-being:

\begin{tcolorbox}[colback=blue!5!white, colframe=blue!75!black]
\textbf{Key Finding:} Creative destruction has heterogeneous welfare effects:
\begin{itemize}
    \item \textbf{Positive for entrants}: New opportunities from innovation
    \item \textbf{Negative for displaced workers}: Short-term unemployment and skill obsolescence
    \item \textbf{Net positive aggregate}: Long-run gains exceed short-run losses
\end{itemize}
\end{tcolorbox}

\subsection{Missing Growth}

Aghion et al. (2019) show that creative destruction leads to measurement problems:

\begin{axiom}[UNMAPPED_PA-4: Missing Growth]
\label{ax:pa4}
Standard GDP measurement misses 0.5-1.0 percentage points of annual growth due to:
\begin{enumerate}
    \item New product varieties not captured by fixed baskets
    \item Quality improvements in existing products
    \item Consumer surplus from creative destruction
\end{enumerate}
\textbf{EBF Connection:} Implies systematic underestimation of welfare gains from innovation.
\end{axiom}

% ============================================================================
% INSTITUTIONS AND REGULATION
% ============================================================================
\section{Institutions, Trust, and Regulation}

\subsection{The Regulation-Distrust Nexus}

Aghion et al. (2010) discover a fundamental tradeoff:

\begin{theorem}[Regulation-Distrust Equilibrium]
Societies can end up in two stable equilibria:
\begin{enumerate}
    \item \textbf{High trust, low regulation}: Trust substitutes for formal regulation
    \item \textbf{Low trust, high regulation}: Distrust necessitates extensive regulation, which further erodes trust
\end{enumerate}
\end{theorem}

\begin{axiom}[UNMAPPED_PA-5: Trust-Regulation Complementarity]
\label{ax:pa5}
Trust ($T$) and regulation ($R$) exhibit negative complementarity:
\begin{equation}
\frac{\partial^2 Output}{\partial T \partial R} < 0
\end{equation}
High trust makes regulation less necessary; extensive regulation crowds out trust.
\textbf{EBF Connection:} Crucial for CONTEXT-CULTURE institutional parameters.
\end{axiom}

\subsection{Distance to Frontier}

Acemoglu, Aghion \& Zilibotti (2006) show that optimal institutions depend on development stage:

\begin{tcolorbox}[colback=green!5!white, colframe=green!75!black]
\textbf{Distance to Frontier Theory:}
\begin{itemize}
    \item \textbf{Far from frontier}: Imitation-based growth; benefit from directed credit, industrial policy
    \item \textbf{Near frontier}: Innovation-based growth; require competitive markets, flexible institutions
\end{itemize}
The ``middle-income trap'' occurs when countries fail to switch institutional modes.
\end{tcolorbox}

% ============================================================================
% GREEN INNOVATION
% ============================================================================
\section{Directed Technical Change and Green Innovation}

\subsection{Path Dependency in Innovation}

Aghion et al. (2016) study how carbon taxes affect innovation direction:

\begin{axiom}[UNMAPPED_PA-6: Directed Technical Change]
\label{ax:pa6}
Innovation direction responds to relative prices with elasticity $\epsilon$:
\begin{equation}
\frac{d \ln(Green\_Patents)}{d \ln(Carbon\_Tax)} = \epsilon \approx 0.3-0.5
\end{equation}
Past innovation in dirty technologies creates path dependency that carbon pricing can redirect.
\textbf{EBF Connection:} Provides foundation for DOMAIN-SUSTAINABILITY interventions.
\end{axiom}

\subsection{Climate Policy Implications}

\begin{tcolorbox}[colback=yellow!10!white, colframe=yellow!75!black, title=Policy Implications]
\begin{enumerate}
    \item Early carbon pricing has outsized effects due to path dependency
    \item R\&D subsidies for green technology complement carbon taxes
    \item Delay is costly: dirty technology lock-in becomes harder to reverse
    \item Optimal policy combines pricing (Pigouvian) with innovation support (Schumpeterian)
\end{enumerate}
\end{tcolorbox}

% ============================================================================
% AI AND AUTOMATION
% ============================================================================
\section{Artificial Intelligence and the Future of Growth}

\subsection{AI as General Purpose Technology}

Aghion, Jones \& Jones (2018) analyze AI's potential impact on growth:

\begin{tcolorbox}[colback=blue!5!white, colframe=blue!75!black]
\textbf{Three Scenarios for AI and Growth:}
\begin{enumerate}
    \item \textbf{Baumol effect}: AI automates routine tasks but bottlenecks in human creativity persist $\Rightarrow$ growth eventually stagnates
    \item \textbf{Steady growth}: AI augments human capabilities proportionally $\Rightarrow$ constant growth rate
    \item \textbf{Singularity}: AI improves AI development $\Rightarrow$ explosive growth
\end{enumerate}
Current evidence suggests Scenario 2 is most likely, but Scenario 1 risks are real.
\end{tcolorbox}

\subsection{Automation and Labor}

Aghion et al. (2020) provide a Schumpeterian perspective on automation:

\begin{itemize}
    \item \textbf{Displacement effect}: Automation replaces workers in specific tasks
    \item \textbf{Productivity effect}: Higher productivity increases demand for complementary labor
    \item \textbf{New task creation}: Innovation creates entirely new job categories
\end{itemize}

Net effect depends on the relative strength of these forces and the speed of adjustment.

% ============================================================================
% EBF INTEGRATION
% ============================================================================
\section{EBF Integration}

\subsection{10C Mapping}

\begin{table}[h]
\centering
\begin{tabular}{lll}
\textbf{10C Element} & \textbf{Aghion Contribution} & \textbf{Parameter} \\
\hline
CORE-WHO (AAA) & Innovators vs. incumbents & Role heterogeneity \\
CORE-WHAT (C) & Welfare from creative destruction & $W_{CD}$ \\
CORE-HOW (B) & Competition-innovation complementarity & Inverted-U \\
CORE-WHEN (V) & Path dependency in innovation & $\delta_{path}$ \\
CORE-WHERE (BBB) & Growth and innovation parameters & $\lambda, \gamma$ \\
CORE-AWARE (AU) & Awareness of technological change & Displacement effects \\
CORE-READY (AV) & Institutional readiness & Distance to frontier \\
CORE-STAGE (AW) & Development stage transitions & Imitation $\rightarrow$ Innovation \\
\end{tabular}
\caption{Aghion contributions mapped to 10C framework}
\end{table}

\subsection{Key Parameter Values for BBB}

\begin{tcolorbox}[colback=green!5!white, colframe=green!75!black, title=Aghion Parameters for BBB Repository]
\begin{itemize}
    \item $\lambda$: 0.04-0.06 - Annual creative destruction rate
    \item $\gamma$: 1.02-1.05 - Quality step size
    \item $HHI^*$: ~0.15 - Optimal competition for innovation
    \item $\epsilon_{carbon}$: 0.3-0.5 - Green innovation elasticity
    \item $\rho_{trust-reg}$: -0.4 - Trust-regulation correlation
    \item Missing growth: 0.5-1.0 pp/year
\end{itemize}
\end{tcolorbox}

% ============================================================================
% WORKED EXAMPLE
% ============================================================================
\section{Worked Example: Optimal Competition Policy}

Consider two industries:
\begin{itemize}
    \item \textbf{Industry A}: Neck-and-neck (HHI = 0.25, tech gap small)
    \item \textbf{Industry B}: Leader-follower (HHI = 0.25, tech gap large)
\end{itemize}

\textbf{Analysis using Aghion framework:}

\textit{Industry A (neck-and-neck):}
\begin{itemize}
    \item Escape competition effect dominates
    \item Increasing competition $\Rightarrow$ more innovation
    \item Policy: Reduce barriers to entry, promote competition
\end{itemize}

\textit{Industry B (leader-follower):}
\begin{itemize}
    \item Schumpeterian effect dominates
    \item Increasing competition $\Rightarrow$ less innovation (reduced rents)
    \item Policy: Allow some market power, support R\&D
\end{itemize}

\textbf{Conclusion:} Uniform competition policy is suboptimal. Context-dependent policy yields higher aggregate innovation.

% ============================================================================
% SUMMARY
% ============================================================================


%%%%%%%%%%%%%%%%%%%%%%%%%%%%%%%%%%%%%%%%%%%%%%%%%%%%%%%%%%%%%%%%%%%%%%%%%%%%%%%
\section{Key Results}
\label{sec:pa-results}
%%%%%%%%%%%%%%%%%%%%%%%%%%%%%%%%%%%%%%%%%%%%%%%%%%%%%%%%%%%%%%%%%%%%%%%%%%%%%%%

The analysis yields the following key results:

\begin{enumerate}
    \item \textbf{Result 1:} Primary finding with quantitative support
    \item \textbf{Result 2:} Secondary finding with implications
    \item \textbf{Result 3:} Practical application or boundary condition
\end{enumerate}


% === AUTO-GENERATED ENTRIES (2026-02-22) ===

%%%%%%%%%%%%%%%%%%%%%%%%%%%%%%%%%%%%%%%%%%%%%%%%%%%%%%%%%%%%%%%%%%%%%%%%%%%%%%%
\subsubsection{1992: An Incomplete Contracts Approach to Financial Contracting}
%%%%%%%%%%%%%%%%%%%%%%%%%%%%%%%%%%%%%%%%%%%%%%%%%%%%%%%%%%%%%%%%%%%%%%%%%%%%%%%

\paragraph{Citation.}
Aghion, Philippe and Bolton, Patrick (1992). ``An Incomplete Contracts Approach to Financial Contracting.'' \textit{Review of Economic Studies}, 59.

\paragraph{10C Integration.}
\begin{itemize}[nosep]
  \item \textbf{Domain}: [To be specified]
  \item \textbf{use\_for}: LIT-AGHION
\end{itemize}


%%%%%%%%%%%%%%%%%%%%%%%%%%%%%%%%%%%%%%%%%%%%%%%%%%%%%%%%%%%%%%%%%%%%%%%%%%%%%%%
\subsubsection{1998: Endogenous Growth Theory}
%%%%%%%%%%%%%%%%%%%%%%%%%%%%%%%%%%%%%%%%%%%%%%%%%%%%%%%%%%%%%%%%%%%%%%%%%%%%%%%

\paragraph{Citation.}
Aghion, Philippe and Howitt, Peter (1998). \textit{Endogenous Growth Theory}. Publisher.

\paragraph{10C Integration.}
\begin{itemize}[nosep]
  \item \textbf{Domain}: [To be specified]
  \item \textbf{use\_for}: LIT-AGHION
\end{itemize}


%%%%%%%%%%%%%%%%%%%%%%%%%%%%%%%%%%%%%%%%%%%%%%%%%%%%%%%%%%%%%%%%%%%%%%%%%%%%%%%
\subsubsection{2005: The Effect of Financial Development on Convergence: Theory a...}
%%%%%%%%%%%%%%%%%%%%%%%%%%%%%%%%%%%%%%%%%%%%%%%%%%%%%%%%%%%%%%%%%%%%%%%%%%%%%%%

\paragraph{Citation.}
Aghion, Philippe and Howitt, Peter and Mayer-Foulkes, David (2005). ``The Effect of Financial Development on Convergence: Theory and Evidence.'' \textit{Quarterly Journal of Economics}, 120.

\paragraph{10C Integration.}
\begin{itemize}[nosep]
  \item \textbf{Domain}: [To be specified]
  \item \textbf{use\_for}: LIT-AGHION
\end{itemize}


%%%%%%%%%%%%%%%%%%%%%%%%%%%%%%%%%%%%%%%%%%%%%%%%%%%%%%%%%%%%%%%%%%%%%%%%%%%%%%%
\subsubsection{2005: Handbook of Economic Growth}
%%%%%%%%%%%%%%%%%%%%%%%%%%%%%%%%%%%%%%%%%%%%%%%%%%%%%%%%%%%%%%%%%%%%%%%%%%%%%%%

\paragraph{Citation.}
Aghion, Philippe and Durlauf, Steven (2005). ``Handbook of Economic Growth.'' \textit{null}, null.

\paragraph{10C Integration.}
\begin{itemize}[nosep]
  \item \textbf{Domain}: [To be specified]
  \item \textbf{use\_for}: LIT-AGHION
\end{itemize}


%%%%%%%%%%%%%%%%%%%%%%%%%%%%%%%%%%%%%%%%%%%%%%%%%%%%%%%%%%%%%%%%%%%%%%%%%%%%%%%
\subsubsection{2010: Volatility and Growth: Credit Constraints and the Compositio...}
%%%%%%%%%%%%%%%%%%%%%%%%%%%%%%%%%%%%%%%%%%%%%%%%%%%%%%%%%%%%%%%%%%%%%%%%%%%%%%%

\paragraph{Citation.}
Aghion, Philippe and Angeletos, George-Marios and Banerjee, Abhijit and Manova, Kalina (2010). ``Volatility and Growth: Credit Constraints and the Composition of Investment.'' \textit{Journal of Monetary Economics}, 57.

\paragraph{10C Integration.}
\begin{itemize}[nosep]
  \item \textbf{Domain}: [To be specified]
  \item \textbf{use\_for}: LIT-AGHION, LIT-O
\end{itemize}


%%%%%%%%%%%%%%%%%%%%%%%%%%%%%%%%%%%%%%%%%%%%%%%%%%%%%%%%%%%%%%%%%%%%%%%%%%%%%%%
\subsubsection{2009: The Causal Impact of Education on Economic Growth: Evidence ...}
%%%%%%%%%%%%%%%%%%%%%%%%%%%%%%%%%%%%%%%%%%%%%%%%%%%%%%%%%%%%%%%%%%%%%%%%%%%%%%%

\paragraph{Citation.}
Aghion, Philippe and Boustan, Leah and Hoxby, Caroline and Vandenbussche, Jerome (2009). ``The Causal Impact of Education on Economic Growth: Evidence from U.S..'' \textit{Working Paper}.

\paragraph{10C Integration.}
\begin{itemize}[nosep]
  \item \textbf{Domain}: [To be specified]
  \item \textbf{use\_for}: LIT-AGHION
\end{itemize}


%%%%%%%%%%%%%%%%%%%%%%%%%%%%%%%%%%%%%%%%%%%%%%%%%%%%%%%%%%%%%%%%%%%%%%%%%%%%%%%
\subsubsection{2008: Academic Freedom, Private-Sector Focus, and the Process of I...}
%%%%%%%%%%%%%%%%%%%%%%%%%%%%%%%%%%%%%%%%%%%%%%%%%%%%%%%%%%%%%%%%%%%%%%%%%%%%%%%

\paragraph{Citation.}
Aghion, Philippe and Dewatripont, Mathias and Stein, Jeremy (2008). ``Academic Freedom, Private-Sector Focus, and the Process of Innovation.'' \textit{RAND Journal of Economics}, 39.

\paragraph{10C Integration.}
\begin{itemize}[nosep]
  \item \textbf{Domain}: [To be specified]
  \item \textbf{use\_for}: LIT-AGHION
\end{itemize}


%%%%%%%%%%%%%%%%%%%%%%%%%%%%%%%%%%%%%%%%%%%%%%%%%%%%%%%%%%%%%%%%%%%%%%%%%%%%%%%
\subsubsection{2009: The Economics of Growth}
%%%%%%%%%%%%%%%%%%%%%%%%%%%%%%%%%%%%%%%%%%%%%%%%%%%%%%%%%%%%%%%%%%%%%%%%%%%%%%%

\paragraph{Citation.}
Aghion, Philippe and Howitt, Peter (2009). \textit{The Economics of Growth}. Publisher.

\paragraph{10C Integration.}
\begin{itemize}[nosep]
  \item \textbf{Domain}: [To be specified]
  \item \textbf{use\_for}: LIT-AGHION
\end{itemize}


%%%%%%%%%%%%%%%%%%%%%%%%%%%%%%%%%%%%%%%%%%%%%%%%%%%%%%%%%%%%%%%%%%%%%%%%%%%%%%%
\subsubsection{2004: Endogenous Political Institutions}
%%%%%%%%%%%%%%%%%%%%%%%%%%%%%%%%%%%%%%%%%%%%%%%%%%%%%%%%%%%%%%%%%%%%%%%%%%%%%%%

\paragraph{Citation.}
Aghion, Philippe and Alesina, Alberto and Trebbi, Francesco (2004). ``Endogenous Political Institutions.'' \textit{Quarterly Journal of Economics}, 119.

\paragraph{10C Integration.}
\begin{itemize}[nosep]
  \item \textbf{Domain}: [To be specified]
  \item \textbf{use\_for}: LIT-AGHION
\end{itemize}


%%%%%%%%%%%%%%%%%%%%%%%%%%%%%%%%%%%%%%%%%%%%%%%%%%%%%%%%%%%%%%%%%%%%%%%%%%%%%%%
\subsubsection{2015: Industrial Policy and Competition}
%%%%%%%%%%%%%%%%%%%%%%%%%%%%%%%%%%%%%%%%%%%%%%%%%%%%%%%%%%%%%%%%%%%%%%%%%%%%%%%

\paragraph{Citation.}
Aghion, Philippe and Dewatripont, Mathias and Du, Luosha and Harrison, Ann and Legros, Patrick (2015). ``Industrial Policy and Competition.'' \textit{American Economic Journal: Macroeconomics}, 7.

\paragraph{10C Integration.}
\begin{itemize}[nosep]
  \item \textbf{Domain}: [To be specified]
  \item \textbf{use\_for}: LIT-AGHION
\end{itemize}


%%%%%%%%%%%%%%%%%%%%%%%%%%%%%%%%%%%%%%%%%%%%%%%%%%%%%%%%%%%%%%%%%%%%%%%%%%%%%%%
\subsubsection{2017: The Social Origins of Inventors}
%%%%%%%%%%%%%%%%%%%%%%%%%%%%%%%%%%%%%%%%%%%%%%%%%%%%%%%%%%%%%%%%%%%%%%%%%%%%%%%

\paragraph{Citation.}
Aghion, Philippe and Akcigit, Ufuk and Hyytinen, Ari and Toivanen, Otto (2017). ``The Social Origins of Inventors.'' \textit{NBER Working Paper}.

\paragraph{10C Integration.}
\begin{itemize}[nosep]
  \item \textbf{Domain}: [To be specified]
  \item \textbf{use\_for}: LIT-AGHION
\end{itemize}


%%%%%%%%%%%%%%%%%%%%%%%%%%%%%%%%%%%%%%%%%%%%%%%%%%%%%%%%%%%%%%%%%%%%%%%%%%%%%%%
\subsubsection{2018: The Impact of Exports on Innovation: Theory and Evidence}
%%%%%%%%%%%%%%%%%%%%%%%%%%%%%%%%%%%%%%%%%%%%%%%%%%%%%%%%%%%%%%%%%%%%%%%%%%%%%%%

\paragraph{Citation.}
Aghion, Philippe and Bergeaud, Antonin and Lequien, Matthieu and Melitz, Marc (2018). ``The Impact of Exports on Innovation: Theory and Evidence.'' \textit{NBER Working Paper}.

\paragraph{10C Integration.}
\begin{itemize}[nosep]
  \item \textbf{Domain}: [To be specified]
  \item \textbf{use\_for}: LIT-AGHION
\end{itemize}


%%%%%%%%%%%%%%%%%%%%%%%%%%%%%%%%%%%%%%%%%%%%%%%%%%%%%%%%%%%%%%%%%%%%%%%%%%%%%%%
\subsubsection{2019: A Theory of Falling Growth and Rising Rents}
%%%%%%%%%%%%%%%%%%%%%%%%%%%%%%%%%%%%%%%%%%%%%%%%%%%%%%%%%%%%%%%%%%%%%%%%%%%%%%%

\paragraph{Citation.}
Aghion, Philippe and Bergeaud, Antonin and Boppart, Timo and Klenow, Peter and Li, Huiyu (2019). ``A Theory of Falling Growth and Rising Rents.'' \textit{NBER Working Paper}.

\paragraph{10C Integration.}
\begin{itemize}[nosep]
  \item \textbf{Domain}: [To be specified]
  \item \textbf{use\_for}: LIT-AGHION
\end{itemize}


%%%%%%%%%%%%%%%%%%%%%%%%%%%%%%%%%%%%%%%%%%%%%%%%%%%%%%%%%%%%%%%%%%%%%%%%%%%%%%%
\subsubsection{2020: The Effects of Automation on Labor: A View from the Schumpet...}
%%%%%%%%%%%%%%%%%%%%%%%%%%%%%%%%%%%%%%%%%%%%%%%%%%%%%%%%%%%%%%%%%%%%%%%%%%%%%%%

\paragraph{Citation.}
Aghion, Philippe and Antonin, Céline and Bunel, Simon (2020). ``The Effects of Automation on Labor: A View from the Schumpeterian Perspective.'' \textit{Working Paper}.

\paragraph{10C Integration.}
\begin{itemize}[nosep]
  \item \textbf{Domain}: [To be specified]
  \item \textbf{use\_for}: LIT-AGHION
\end{itemize}


%%%%%%%%%%%%%%%%%%%%%%%%%%%%%%%%%%%%%%%%%%%%%%%%%%%%%%%%%%%%%%%%%%%%%%%%%%%%%%%
\subsubsection{2015: Lessons from Schumpeterian Growth Theory}
%%%%%%%%%%%%%%%%%%%%%%%%%%%%%%%%%%%%%%%%%%%%%%%%%%%%%%%%%%%%%%%%%%%%%%%%%%%%%%%

\paragraph{Citation.}
Aghion, Philippe and Akcigit, Ufuk and Howitt, Peter (2015). ``Lessons from Schumpeterian Growth Theory.'' \textit{American Economic Review}, 105.

\paragraph{10C Integration.}
\begin{itemize}[nosep]
  \item \textbf{Domain}: [To be specified]
  \item \textbf{use\_for}: LIT-AGHION
\end{itemize}


%%%%%%%%%%%%%%%%%%%%%%%%%%%%%%%%%%%%%%%%%%%%%%%%%%%%%%%%%%%%%%%%%%%%%%%%%%%%%%%
\subsubsection{2017: Tax Simplicity and Heterogeneous Learning}
%%%%%%%%%%%%%%%%%%%%%%%%%%%%%%%%%%%%%%%%%%%%%%%%%%%%%%%%%%%%%%%%%%%%%%%%%%%%%%%

\paragraph{Citation.}
Aghion, Philippe and Akcigit, Ufuk and Lequien, Matthieu and Stantcheva, Stefanie (2017). ``Tax Simplicity and Heterogeneous Learning.'' \textit{NBER Working Paper}.

\paragraph{10C Integration.}
\begin{itemize}[nosep]
  \item \textbf{Domain}: [To be specified]
  \item \textbf{use\_for}: LIT-AGHION
\end{itemize}


%%%%%%%%%%%%%%%%%%%%%%%%%%%%%%%%%%%%%%%%%%%%%%%%%%%%%%%%%%%%%%%%%%%%%%%%%%%%%%%
\subsubsection{2023: Environmental Preferences and Technological Choices: Is Mark...}
%%%%%%%%%%%%%%%%%%%%%%%%%%%%%%%%%%%%%%%%%%%%%%%%%%%%%%%%%%%%%%%%%%%%%%%%%%%%%%%

\paragraph{Citation.}
Aghion, Philippe and Bénabou, Roland and Martin, Ralf and Roulet, Alexandra (2023). ``Environmental Preferences and Technological Choices: Is Market Competition Clean or Dirty?.'' \textit{American Economic Review: Insights}, 5.

\paragraph{10C Integration.}
\begin{itemize}[nosep]
  \item \textbf{Domain}: [To be specified]
  \item \textbf{use\_for}: LIT-AGHION
\end{itemize}


%%%%%%%%%%%%%%%%%%%%%%%%%%%%%%%%%%%%%%%%%%%%%%%%%%%%%%%%%%%%%%%%%%%%%%%%%%%%%%%
\subsubsection{2024: The Economics of AI}
%%%%%%%%%%%%%%%%%%%%%%%%%%%%%%%%%%%%%%%%%%%%%%%%%%%%%%%%%%%%%%%%%%%%%%%%%%%%%%%

\paragraph{Citation.}
Aghion, Philippe and Jones, Benjamin and Jones, Charles (2024). ``The Economics of AI.'' \textit{Annual Review of Economics}, 16.

\paragraph{10C Integration.}
\begin{itemize}[nosep]
  \item \textbf{Domain}: [To be specified]
  \item \textbf{use\_for}: LIT-AGHION
\end{itemize}


%%%%%%%%%%%%%%%%%%%%%%%%%%%%%%%%%%%%%%%%%%%%%%%%%%%%%%%%%%%%%%%%%%%%%%%%%%%%%%%
\subsubsection{2021: Decentralization, Economic Shocks, and Organization}
%%%%%%%%%%%%%%%%%%%%%%%%%%%%%%%%%%%%%%%%%%%%%%%%%%%%%%%%%%%%%%%%%%%%%%%%%%%%%%%

\paragraph{Citation.}
Philippe Aghion and Nick Bloom and Raffaella Sadun and John Van Reenen (2021). ``Decentralization, Economic Shocks, and Organization.'' \textit{American Economic Journal: Applied Economics}, 13.

\paragraph{10C Integration.}
\begin{itemize}[nosep]
  \item \textbf{Domain}: [To be specified]
  \item \textbf{use\_for}: LIT-AGHION
\end{itemize}

% === END AUTO-GENERATED ENTRIES ===
\section{Summary}

\begin{tcolorbox}[colback=gray!10!white, colframe=gray!75!black, title=Key Takeaways]
\begin{enumerate}
    \item \textbf{Schumpeterian Growth:} Innovation and creative destruction drive long-run prosperity
    \item \textbf{Inverted-U:} Competition-innovation relationship is non-monotonic
    \item \textbf{Missing Growth:} Creative destruction causes systematic underestimation of welfare gains
    \item \textbf{Trust-Regulation:} Societies face a fundamental institutional tradeoff
    \item \textbf{Green Innovation:} Carbon pricing redirects technical change with path dependency
    \item \textbf{AI and Growth:} Automation effects depend on new task creation and adjustment speed
\end{enumerate}
\end{tcolorbox}

% ============================================================================
% GLOSSARY
% ============================================================================
\section{Glossary}

See Appendix G (Glossary) for standard EBF terms. Aghion-specific terms:

\begin{description}
    \item[Creative Destruction] Process by which new innovations render existing technologies obsolete
    \item[Distance to Frontier] Gap between a country's productivity and the global technology frontier
    \item[Escape Competition Effect] Incentive for neck-and-neck firms to innovate to gain market advantage
    \item[Inverted-U] Non-monotonic relationship where intermediate values optimize outcomes
    \item[Schumpeterian Effect] Reduced innovation incentive when competition erodes post-innovation rents
\end{description}

% ============================================================================
% CRITICAL FOUNDATIONS
% ============================================================================
\section{Critical Foundations}

\begin{enumerate}
    \item Aghion, P. \& Howitt, P. (1992). A Model of Growth Through Creative Destruction. \textit{Econometrica}. [Foundational]
    \item Aghion, P. et al. (2005). Competition and Innovation: An Inverted-U Relationship. \textit{QJE}. [Competition]
    \item Aghion, P. et al. (2010). Regulation and Distrust. \textit{QJE}. [Institutions]
    \item Aghion, P. et al. (2016). Carbon Taxes, Path Dependency, and Directed Technical Change. \textit{JPE}. [Climate]
    \item Aghion, P., Jones, B. \& Jones, C. (2018). Artificial Intelligence and Economic Growth. [AI]
\end{enumerate}

% ============================================================================
% OPEN ISSUES
% ============================================================================
\section{Open Issues}

\begin{enumerate}
    \item \textbf{Measuring creative destruction:} How to better quantify innovation obsolescence?
    \item \textbf{Distributional effects:} How to compensate losers from creative destruction?
    \item \textbf{AI trajectory:} Which growth scenario will materialize?
    \item \textbf{Global coordination:} How to coordinate innovation policy across countries?
\end{enumerate}

% ============================================================================
% REFERENCES
% ============================================================================
\section{References}

\nocite{bcm_master}

For complete bibliography, see Master Bibliography (\texttt{bibliography/master\_references.bib}).

Key Aghion references integrated in this appendix:
\begin{itemize}[nosep]
    \item aghion1992model, aghion2005competition, aghion2001competition
    \item aghion2010regulation, aghion2016carbon, aghion2019missing
    \item aghion2018artificial, aghion1997formal, aghion2021power
    \item aghion2014what, aghion2016creative
\end{itemize}

\end{document}
